\section{A bounded check of Street's Baroco reconstruction}
\label{sec:annex-p13-street-arabic-2004}

\begin{annexitem}{annex:p13-street2004-baroco}{thm:p13-street-2004-baroco}{The two-stage classical argument at fixed interpretation}
\begin{proof}
Let \(C\) abbreviate \(\exists x(J(x)\land\neg B(x))\).
Under classical first-order semantics, \(\neg C\) means
that every \(J\)-individual is a \(B\)-individual. With
\(\forall x(B(x)\to D(x))\), each such individual is
therefore a \(D\)-individual. This establishes the
conditional from \(\neg C\) to
\(\forall x(J(x)\to D(x))\), keeping the first premise fixed.

The second premise supplies an individual which is \(J\)
and not \(D\), so that universal consequent is false.
Consequently \(\neg C\) is false, and classical double
negation gives \(C\). Equivalently, choose the second
premise's witness: it cannot be \(B\), because the first
premise would then make it \(D\). This also verifies
the conclusion without an extra existence premise for \(B\).

The source's conjunctive/exceptive decomposition is thus
checked in the declared fixed classical interpretation.
The proof does not identify its ordinary implication or
negation with every historical quantified conditional,
temporal absolute or modal proposition.
\end{proof}
\end{annexitem}
